\chapter{Testing and validation}
The testing and validation phase was carried out to verify that the application worked as a functional Android prototype and that the experience was understandable and satisfactory for potential users. Since \textit{Micorrizas} combines mobile interaction, geolocation, multiplayer synchronisation and outdoor exploration, validation is complex, as we need to take x users (between 2 and 6) to the \textit{Jardín de los Sentidos}. Furthermore, each user must have an Android device with geolocation capability.

First, the application was tested on Android devices in order to check the main gameplay flow: creating a room, joining a cooperative session, accessing the map, detecting the player’s position, activating missions, completing challenges and displaying the final result. Special attention was paid to the geolocation system, as this was the most challenging feature to get up and running; furthermore, the game relies on the player’s physical presence in the \textit{Jardín de los Sentidos}, as this is where the teaching activity takes place. For this reason, several tests had to be carried out in the real-world environment, as the behaviour of the GPS could not be reliably validated within the Unity editor.

In addition to technical testing, the game was tested with a group of users in order to collect feedback about the clarity of the interface, the understanding of the missions and the overall flow of the experience. During these tests, a think-aloud protocol was used. Participants were encouraged to verbalise what they were thinking while interacting with the application, which made it possible to identify confusing instructions, unclear visual elements or moments where the game flow was not sufficiently intuitive.

The first round of feedback focused mainly on the structure of the missions and the behaviour of the map. As a result of this feedback, images were added to the second mission in order to support the sound-recognition task and make the available options easier to understand. The activation areas on the map were also delimited more clearly, so that players could better understand where they needed to be in order to start a mission. What’s more, when you complete the mission, there’s visual feedback, the colour changes to indicate that that area of the garden has now been restored.

Other suggestions from this first round were considered but not finally implemented. One of them was the addition of a separate learning scene before the second mission. This idea was rejected because it interrupted the rhythm of the experience and added an extra step before the mission itself. Another suggestion was to prevent players from removing pairs in the fourth mission and to add an additional learning phase before it. This was also not implemented in the final version, as it would have increased the complexity of the mission and moved it away from the intended short and direct interaction model.

A second round of feedback was provided by the lecturers involved in the original educational activity. This feedback was especially relevant because it helped to adjust the prototype to the pedagogical context of the subject. In the fourth quest, an introductory reference to flavours was added so that players could better understand the relationship between the activity and the Garden of Taste. In the first quest, the plant images were enlarged to make visual recognition easier on mobile screens. In the fourth mission, an additional pressure element was introduced to make the activity more dynamic. The user interface was redistributed in some scenes to improve readability and interaction, and the avatar of the GPS indicator was adjusted so that players could more easily interpret the location state of the application.

The validation process helped to refine the game from both a technical and an educational point of view. The tests confirmed that the application could run as an Android prototype and that the main cooperative flow was functional. At the same time, the feedback received made it possible to improve the clarity of the missions, the readability of the interface and the connection between the game mechanics and the educational activity on which the project is based.
